\documentclass[a4paper,12pt]{article}
\usepackage[utf8]{inputenc}
\usepackage[T1]{fontenc}
\usepackage[ngerman]{babel}
\usepackage{geometry}
\usepackage{booktabs}
\usepackage{amsmath}
\geometry{margin=2.5cm}

\begin{document}

\section*{Fourier Neural Operator Based Fluid--Structure Interaction for
Predicting the Vesicle Dynamics}

\textbf{Autoren:} Wang Xiao, Ting Gao, Kai Liu, Jinqiao Duan, Meng Zhao\\
\textbf{Jahr:} 2024 \quad \textbf{Journal:} arXiv:2401.02311 (eingereicht bei Elsevier)

% -------------------------------------------------------------------
\subsection*{Motivation und Problemstellung}
% -------------------------------------------------------------------

Fluid-Struktur-Interaktionsprobleme (FSI) -- bei denen ein deformierbares
Objekt in eine viskose Strömung eingebettet ist -- erfordern typischerweise
aufwendige numerische Verfahren wie die \emph{Immersed Boundary Method}~(IBM).
Diese klassischen CFD-Löser sind für große Parameterräume oder lange
Simulationszeiträume rechnerisch prohibitiv teuer, da bei jedem Zeitschritt
ein vollständiges Strömungsfeld neu berechnet werden muss.

Die Arbeit begegnet diesem Engpass, indem der Fourier Neural Operator~(FNO)
als datengetriebener Strömungslöser eingesetzt wird, der anschließend mit
einem klassischen Strukturlöser (Finite-Differenzen-Methode) gekoppelt wird.
Das Anwendungsbeispiel ist die Dynamik von \emph{Vesikeln} -- membranumhüllten
Flüssigkeitsbläschen -- in einer zeitabhängigen Dehnströmung, die durch die
inkompressible Stokes-Gleichung beschrieben wird:
\begin{align}
    \rho \frac{\partial \mathbf{u}^{\text{ind}}}{\partial t}
    &= \mu \,\Delta \mathbf{u}^{\text{ind}} - \nabla p + \mathbf{f}, \\
    \nabla \cdot \mathbf{u}^{\text{ind}} &= 0.
\end{align}
Die Kopplung zwischen Fluid und Membran erfolgt über die Dirac-Delta-Funktion
der IBM.  Der FNO übernimmt dabei ausschließlich die zeitliche Propagation des
Strömungsfeldes; die Strukturgleichungen werden weiterhin klassisch gelöst.

% -------------------------------------------------------------------
\subsection*{Methodik: FNO-basierter FSI-Löser}
% -------------------------------------------------------------------

Der FNO lernt eine Abbildung zwischen Funktionenräumen: gegeben $\alpha$ aufeinanderfolgende Zeitschritte des Geschwindigkeitsfeldes wird der nächste
Zeitschritt prädiziert. Die Architektur besteht aus einem
\emph{Lifting-Layer}~$P$, mehreren Fourier-Integral-Operator-Schichten
\begin{equation}
    v_{t+1}(x) = \sigma\!\left(W v_t(x)
    + \mathcal{F}^{-1}\!\left[R \cdot \mathcal{F}[v_t]\right](x)\right),
\end{equation}
sowie einem abschließenden \emph{Projection-Layer}~$Q$.  Die komplexwertigen
Gewichtsmatrizen~$R$ werden direkt als lernbare Fourier-Koeffizienten
parametrisiert und auf $k_{\max} = 12$ Moden trunkiert.

Trainiert wird auf 3D-Simulationsdaten ($64^3$-Gitter) der IBM für zwei
verschiedene Anfangsformen des Vesikels.  Das Zeitfenster beträgt $\alpha = 10$
Eingabeschritte.  Es werden vier Experimente (FNO-1 bis FNO-4) systematisch
verglichen, die sich in zwei Dimensionen unterscheiden:

\begin{center}
\begin{tabular}{llll}
\toprule
Experiment & Trainingsmethode & Trainingsdaten & Samples \\
\midrule
FNO-1 & sequence-to-one   & ohne Gleichgewichtszustand & 340 \\
FNO-2 & sequence-to-one   & mit Gleichgewichtszustand  & 420 \\
FNO-3 & sequence-to-sequence & ohne Gleichgewichtszustand & 320 \\
FNO-4 & sequence-to-sequence & mit Gleichgewichtszustand  & 400 \\
\bottomrule
\end{tabular}
\end{center}

% -------------------------------------------------------------------
\subsection*{Zentrale Ergebnisse}
% -------------------------------------------------------------------

\paragraph{Sequence-to-one vs.\ sequence-to-sequence.}
Die \emph{sequence-to-one}-Variante (FNO-1, FNO-2) erzielt in den ersten
Zeitschritten geringere Fehler, da keine kumulative Fehlerakkumulation stattfindet. Für Langzeitprognosen (ab dem 15.\ Zeitschritt) erweist sich dagegen
die \emph{sequence-to-sequence}-Variante (FNO-3, FNO-4) als überlegen: Die
gleichzeitige Optimierung über mehrere Zeitschritte reguliert die
Fehlerakkumulation implizit.

\paragraph{Interpolation vs.\ Extrapolation.}
Wie bei anderen neuronalen Operatoren ist die Vorhersagegenauigkeit für
\emph{Interpolation} (Eingabedaten aus dem Trainingsbereich) stets höher
als für \emph{Extrapolation} (Eingabedaten aus einem noch nicht gesehenen
Zeitbereich).  Dies gilt unabhängig davon, ob Gleichgewichtsdaten im
Trainingsset enthalten sind.

\paragraph{Einfluss der Gleichgewichtsdaten.}
Das Einbeziehen des stationären Zustands in die Trainingsdaten verbessert die
Vorhersagegenauigkeit geringfügig, da das Netzwerk die asymptotische Lösung
der Differentialgleichung explizit sieht.

\paragraph{Vesikeldynamik.}
Die aus dem FNO-Strömungsfeld abgeleiteten Vesikelkonfigurationen stimmen
qualitativ gut mit den IBM-Referenzlösungen überein.  Der mittlere absolute
Fehler der Vesikelposition liegt in allen Szenarien unter 1.4\,\% des
Gesamtdomänendurchmessers.  Der FNO-basierte FSI-Löser ist dabei um eine
Größenordnung schneller als der klassische IBM-Löser.

% -------------------------------------------------------------------
\subsection*{Bezug zur Masterarbeit}
% -------------------------------------------------------------------

Diese Arbeit weist mehrere direkte Parallelen zur vorliegenden Masterarbeit
auf und ist in verschiedener Hinsicht relevant:

\paragraph{Gemeinsamer physikalischer Rahmen.}
Beide Arbeiten operieren im Regime inkompressibler Stokes-Strömung mit sehr
kleinen Reynolds-Zahlen und eingebetteten Objekten (Vesikel bzw.\ Kristalle).
Die zugrundeliegenden Differentialgleichungen sind strukturell identisch; die
Randbedingungen und Geometrien unterscheiden sich jedoch: Vesikel sind
deformierbare, membranumhüllte Objekte, während die Kristalle in der
Masterarbeit als rigide Kreise modelliert werden.  Diese Vereinfachung
ermöglicht es der Masterarbeit, sich vollständig auf die Frage der
Architekturwahl und Generalisierung zu konzentrieren, ohne die FSI-Kopplung
zu berücksichtigen.

\paragraph{FNO als gemeinsame Architektur.}
Xiao et al.\ setzen den FNO für zeitabhängige Strömungsvorhersage ein,
während die Masterarbeit den FNO für stationäre Strömungsfelder mit variabler
Kristallanzahl verwendet.  Beide Kontexte testen die Generalisierungsfähigkeit
des FNO auf ungesehene geometrische Konfigurationen -- ein zentrales Thema
beider Arbeiten.  Die in Xiao et al.\ beobachtete globale Erfassung von
Stokes-Interaktionen durch die spektrale Repräsentation des FNO stimmt mit
der theoretischen Motivation in der Masterarbeit überein: Die elliptische
Natur der Stokes-Gleichungen begünstigt Architekturen mit globalem
Rezeptivfeld.

\paragraph{Komplementäre Architekturvergleiche.}
Während Xiao et al.\ ausschließlich den FNO untersuchen, ist ein zentraler
Beitrag der Masterarbeit der \emph{systematische Vergleich} von U-Net und
FNO auf demselben physikalischen Problem.  Die Erkenntnisse aus Xiao et al.\
zu Stärken und Schwächen des FNO (etwa die Anfälligkeit für Extrapolation)
liefern dabei wertvolle Erwartungshaltungen für diesen Vergleich.

\paragraph{Stream-Function-Formulierung.}
In der Masterarbeit wird die inkompressible Strömung über die Stromfunktion
$\psi$ prädiziert, was die Divergenzfreiheit analytisch erzwingt.  Xiao et al.\
prädizieren direkt das Geschwindigkeitsfeld und müssen die
Divergenzfreiheitsbedingung daher implizit durch die Trainingsdaten und den
FNO-Bias gewährleisten.  Die Überlegenheit der Stromfunktionsformulierung
als physikalisch strukturiertes Lernziel ist damit ein methodischer
Mehrwert der Masterarbeit gegenüber Xiao et al.

\paragraph{Statische vs.\ zeitabhängige Prognose.}
Ein wesentlicher Unterschied besteht im Prognosetyp: Xiao et al.\ lösen ein
Zeitentwicklungsproblem, bei dem der FNO autoregressiv eingesetzt wird.  Die
Masterarbeit hingegen betrachtet \emph{stationäre} Strömungsfelder, was das
Problem strukturell vereinfacht (kein Fehlerakkumulationseffekt), jedoch die
Generalisierung auf geometrisch neuartige Konfigurationen in den Vordergrund
stellt.  Die Erkenntnis von Xiao et al., dass sequence-to-sequence-Training
die Langzeitstabilität verbessert, ist für die Masterarbeit zwar nicht direkt
übertragbar, unterstreicht aber grundsätzlich den Wert von Regularisierungsstrategien, die die physikalische Konsistenz der Ausgabe fördern.

\end{document}